% --------------------------------------------------------------
%  Simple LuaLaTeX template – “Zeugnis”
%  Reads data from   student_data.lua   (see sample next to this file)
%  and produces a 2‑page A4 report similar to the HTML version.
% --------------------------------------------------------------

% --- Preamble -------------------------------------------------
% ----------------------------------------------------
%  Minified, LuaLaTeX-friendly preamble for “Zeugnis”
% ----------------------------------------------------
\documentclass[12pt,a4paper]{article}

% ---------------- Page layout & colours -------------
\usepackage[top=2cm,bottom=2cm,
left=1.5cm,right=1.5cm]{geometry}

\usepackage[table]{xcolor}
\definecolor{egeblue}{RGB}{0,99,142}
\definecolor{egelightblue}{RGB}{66,144,179}
\definecolor{greenEnglish}{rgb}{0,0.5,0}

% ---------------- Language --------------------------
\usepackage[german]{babel}      % LuaLaTeX handles UTF-8 natively

% ---------------- Fonts -----------------------------
\usepackage{fontspec}           % modern font loader
\setmainfont{Latin Modern Sans}        % or any installed font
%\setsansfont{Helvetica Neue}            % replaces \usepackage{helvet}
%\setsansfont{Latin Modern Sans}
\usepackage{amsmath,amssymb}
%\usepackage{unicode-math}       % full Unicode maths
%\setmathfont{Latin Modern Math}

% If you need one single sans-serif default:
% \renewcommand*\familydefault{\sfdefault}

% ---------------- Micro-typography ------------------
\usepackage{microtype}          % protrusion & font expansion
\usepackage{selnolig}           % fixes spurious ligatures

% ---------------- Headings, spacing -----------------
\usepackage{titlesec}
\usepackage{setspace}           % \singlespacing, etc.
\usepackage{ragged2e}
\usepackage{enumitem}

% ---------------- Lua helpers -----------------------
\usepackage{luacode}

% ---------------- Graphics / header / footer --------
\usepackage{graphicx}
\usepackage{fancyhdr}

% ---------------- Tables & boxes --------------------
\usepackage{nowidow}            % prevents widows/orphans

\usepackage{tabularray}

% ---------------- Maths & symbols -------------------
\usepackage{amsmath,amssymb}

% ---------------- Strikeout, underline --------------
\usepackage[normalem]{ulem}      % keep \emph italic, enable \sout

% ---------------- Misc helpers ----------------------
\usepackage{xstring,etoolbox}

% If you still need a *specific* Unicode symbol that your fonts
% don’t contain, you can define it once, e.g.:
% \newunicodechar{⟨}{\ensuremath\langle}

% ----------------------------------------------------
\setlength{\parindent}{0pt}
\newcommand{\tableGradeWidth}{.080\linewidth}
\newcommand{\tableGradeWidthLevel}{.06\linewidth}

\newcommand{\schoolReportDay}{27.06.2025}

\makeatletter
\def\nobreakhline{%
	\noalign{\ifnum0=`}\fi
	\penalty\@M
	\futurelet\@let@token\LT@@nobreakhline}
\def\LT@@nobreakhline{%
	\ifx\@let@token\hline
	\global\let\@gtempa\@gobble
	\gdef\LT@sep{\penalty\@M\vskip\doublerulesep}% <-- change here
	\else
	\global\let\@gtempa\@empty
	\gdef\LT@sep{\penalty\@M\vskip-\arrayrulewidth}% <-- change here
	\fi
	\ifnum0=`{\fi}%
	\multispan\LT@cols
	\unskip\leaders\hrule\@height\arrayrulewidth\hfill\cr
	\noalign{\LT@sep}%
	\multispan\LT@cols
	\unskip\leaders\hrule\@height\arrayrulewidth\hfill\cr
	\noalign{\penalty\@M}%
	\@gtempa}
\makeatother


\newcommand\competencyLevel{
	\begin{tabularx}{\linewidth}{|X|C{\tableGradeWidth}|C{\tableGradeWidth}|C{\tableGradeWidth}|C{\tableGradeWidth}|}
		\hline
		\textbf{Kompetenzfelder und Lernziele der Fachbereiche} & \rot{\parbox{2.7cm}{\thead{sehr gut erfüllt}}} & \rot{\thead{gut erfüllt}} & \rot{\thead{teilweise}} & \rot{\thead{nicht erfüllt}} \\
		\hline
	\end{tabularx}
}

\newcommand{\gradeOne}{& $\boxtimes$ & $\Box$ & $\Box$ & $\Box$\\\nobreakhline}
\newcommand{\gradeTwo}{& $\Box$ & $\boxtimes$ & $\Box$ & $\Box$\\\nobreakhline}
\newcommand{\gradeThree}{& $\Box$ & $\Box$ & $\boxtimes$ & $\Box$\\\nobreakhline}
\newcommand{\gradeFour}{& $\Box$ & $\Box$ & $\Box$ & $\boxtimes$\\\nobreakhline}
\newcommand{\gradeNotGiven}{& \multicolumn{4}{c|}{nicht erteilt}\\\nobreakhline}
\newcommand{\gradeComesWithSecondHalfYear}{& \multicolumn{4}{c|}{wird im 2. Halbjahr belegt}\\\nobreakhline}
%\newcommand{\gradeDefault}{&  &  &  & \\\nobreakhline}
\newcommand\gradeDefault{& \multicolumn{4}{c|}{\textemdash} \\}
\newcommand{\gradeSkip}{& $\Box$ & $\Box$ & $\Box$ & $\Box$\\\nobreakhline}


\newcommand\levelRow[1]{\multicolumn{5}{|c|}{#1}\\\hline}
\newcommand{\levelOne}{\levelRow{\centering Du hast vorwiegend auf {\color{greenEnglish} Anforderungsebene 1} gearbeitet.}}
\newcommand{\levelTwo}{\levelRow{\centering Du hast vorwiegend auf {\color{blue} Anforderungsebene 2} gearbeitet.}}
\newcommand{\levelThree}{\levelRow{\centering Du hast vorwiegend auf {\color{red} Anforderungsebene 3} gearbeitet.}}
\newcommand{\levelSeven}{\levelRow{\centering bis Klasse 7 ohne Anforderungsebene}}
\newcommand{\levelEight}{\levelRow{\centering bis Klasse 8 ohne Anforderungsebene}}
\newcommand{\levelNine}{\levelRow{\centering bis Klasse 9 ohne Anforderungsebene}}
\newcommand{\noLevel}{}
\newcommand{\compText}[1]{\multicolumn{5}{|p{0.975\linewidth}|}{#1}\\\hline}


% #1 ... competence
% #2 ... grade
\newcommand\competencyMajorSubject[2]{\makecell[{{>{\parindent0em}p{\linewidth}}}]{#1} #2}

% #1 ... subject
% #2 ... competencies
% #3 ... level
\newcommand\competencyTableMajorSubject[3]{
	\renewcommand{\arraystretch}{1.5}
	\begin{xltabular}{\linewidth}{|K{X}|C{\tableGradeWidth}|C{\tableGradeWidth}|C{\tableGradeWidth}|C{\tableGradeWidth}|}
		\nobreakhline
		\multicolumn{5}{|>{\columncolor{egelightblue}}l|}{\textit{\textbf{#1}}}\\
		\nobreakhline
		#2
		#3\\
	\end{xltabular}\vspace{2pt}
}

\newcommand\competencyTableMajorSubjectWord[2]{
	\renewcommand{\arraystretch}{1.5}
	\begin{xltabular}{\linewidth}{|K{X}|C{\tableGradeWidth}|C{\tableGradeWidth}|C{\tableGradeWidth}|C{\tableGradeWidth}|}
		\nobreakhline
		\multicolumn{5}{|>{\columncolor{egelightblue}}l|}{\textit{\textbf{#1}}}\\
		\nobreakhline
		\multicolumn{5}{|>K{X}|}{#2}\\
	\end{xltabular}\vspace{2pt}
}

% argument is the page threshold until a newline should be added
\newcommand\newOptionalNewPage[1]{
	\ifnum #1>\value{page}
		\newpage
		\thispagestyle{plain}
		\phantom{~}
	\fi
}

\newcommand{\newpagedefs}{
	\newpage
	\newgeometry{headheight=90pt,top=7cm,
		bottom=2cm, 
		inner=2cm,
		outer=2cm}
	\setlength{\headsep}{-0.4cm}
	\pagestyle{fancy}
}

\newcommand{\newpagedefsLastPage}{
	\newpage
	\newgeometry{top=2cm,
		bottom=2cm, 
		inner=2cm,
		outer=2cm}
	\pagestyle{fancy}
}

\newcommand\formatText[1]{
	\noexpandarg
	\StrBefore{#1}{\\}[\studentIntro]
	\StrBehind{#1}{\\}[\certText]
	{\Large\setstretch{1.10}\textbf{\LARGE\studentIntro\vspace{.5em}\\}\nowidow[11]\noclub[11]\certText\par}
}

%%\newcommand*\formatTable[1]{%
%%	\IfSubStr{#1}{\hline}{
%%		\StrCount{#1}{\hline}[\occur]%
%%		\StrBefore[\occur]{#1}{\hline}[\strBefore]
%%		\StrBehind[\occur]{#1}{\hline}[\lastHline]}
%%		{\strBefore} 
%%%		\StrSubstitute{\lastHline}{\hline}{\nobreakhline} 
%%%		\\
%%	}
%%}
%
%\newcommand*\formatTable[1]{%
%	\IfSubStr{#1}{\hline}
%	{
%		\StrCount{#1}{\hline}[\occur]
%		\StrBefore[\occur]{#1}{\hline}
%	}{}
%}

\pagestyle{fancy}

\fancypagestyle{mypagestyle}{%
	\chead{\competencyLevel}
	\cfoot{\small\raggedright\textit{Legende}:	Anforderungsebene Grün = Hauptschule ( AE I ), Anforderungsebene Blau = Realschule ( AE II ), Anforderungsebene Rot = Gymnasium ( AE III ),	n.b. = nicht bewertet}
}

\renewcommand{\headrulewidth}{0pt}

\widowpenalty10000
\clubpenalty10000

\renewcommand\cellalign{tl}  % top-left alignment inside every \makecell


% --- Lua helpers ----------------------------------------------
\begin{luacode*}
	-- Load student data table
	dofile("student_data.lua")
	
	-- Escape TeX‑special characters coming from Lua strings
	 function tex_escape(str)
	 return (str
		:gsub("([%%#$&{}_\\])", "\\%1")   -- escape % # $ & _ { } \
		:gsub("\n", "\\\\"))              -- newline → \\
  	 end
	
	-- Convenience for checkbox columns (1‑4)
	function checkbox(level, col)
	return (level == col) and "$\\boxtimes$" or "$\\Box$"
	end
\end{luacode*}

\newcommand{\headcol}[2]{\cellcolor{#1}\color{white}\bfseries #2}

% --- Document --------------------------------------------------
\begin{document}
	
	% ---------- Title page ----------------------------------------
		\pagestyle{plain}
	\fcolorbox{egeblue}{white}{
		\begin{minipage}[t][.98\textheight][t]{.97\textwidth}
			\begin{center}\vspace{1cm}
				\includegraphics[width=.9\linewidth]{Logotop_mit_ESM_zeugnis}
			\end{center}
			\vspace{3cm}
			\begin{center}\setstretch{1.5}
				\textbf{\Huge Zeugnis}\\
				\textbf{\LARGE ~\\
					Evangelische Gemeinschaftsschule Erfurt\\\vspace{2cm}
					\directlua{tex.sprint(student.part_of_year)} \directlua{tex.sprint(student.school_year)}}
			\end{center}\vspace{3cm}
			\begin{center}\LARGE
				\begin{tabular}{ll}
					\textbf{Klasse:} &  \directlua{tex.sprint(student.classRoom)}\vspace{1cm}\\
					\textbf{Name:} &   \directlua{tex.sprint(student.first_name)} \directlua{tex.sprint(student.last_name)}\\%
					& geb.: \directlua{tex.sprint(student.date_of_birth)}\\%
				\end{tabular}
			\end{center}
			\vfill
			\begin{center}
				\includegraphics[width=0.4\linewidth]{Stiftung}
			\end{center}
			\vspace{1cm}
		\end{minipage}
	}%


	\newpage
	
	% ---------- Personal text -------------------------------------
	Lieber \directlua{tex.sprint(student.first_name)},\\
	\directlua{tex.sprint(student.personal_text)}
	\newpage
	\vspace{1em}
	
%	% ---------- Subjects & competences ----------------------------

%\begin{luacode*}
%	local gradeCmd = {
%		[1]  = "\\gradeOne",
%		[2]  = "\\gradeTwo",
%		[3]  = "\\gradeThree",
%		[4]  = "\\gradeFour",
%		ne   = "\\gradeNotGiven",
%		HJ2  = "\\gradeComesWithSecondHalfYear",
%		default = "\\gradeDefault"
%	}
%
%	local headerTemplate = [=[
%	\cellcolor{egelightblue}{\textit{\textbf{%s}}} &
%		\makecell[c]{\footnotesize{sehr~gut\\erfüllt}} &
%		\makecell[c]{\footnotesize{gut\\erfüllt}} &
%		\makecell[c]{\footnotesize{teilweise\\erfüllt}} &
%		\makecell[c]{\footnotesize{nicht\\erfüllt}}\\
%		]=]
%
%for _, subj in ipairs(student.subjects) do
%	tex.print("\\renewcommand{\\arraystretch}{1.5}")
%	tex.print("\\begin{xltabular}{\\linewidth}{|X|C{\\tableGradeWidth}|C{\\tableGradeWidth}|C{\\tableGradeWidth}|C{\\tableGradeWidth}|}")
%	tex.print("\\nobreakhline")
%	-- tex.print(string.format(headerTemplate, tex_escape(subj.name)))
%	tex.print("\\nobreakhline")
%	
%	-- ----------------------------------------------------------
%	-- 2. exactly one row per AREA  ➔ all competences in 1st cell
%	-- ----------------------------------------------------------
%    for _, area in ipairs(subj.areas) do
%		-- build the first (description) cell with makecell
%		local lines = { "\\textbf{" .. tex_escape(area.title) .. "}" }
%		for _, comp in ipairs(area.competences) do
%			table.insert(lines, tex_escape(comp.description))
%		end
%		local firstCell = "\\makecell[tl]{" .. table.concat(lines, "\\\\") .. "}"
%		
%		-- print the complete row
%		tex.print(firstCell .. " " .. (gradeCmd[area.grade] or "\\default"))
%	end
%	
%	tex.print("\\end{xltabular}")
%end
%\end{luacode*}

\begin{xltabular}{\linewidth}{|X|C{\tableGradeWidth}|C{\tableGradeWidth}|C{\tableGradeWidth}|C{\tableGradeWidth}|}
	\hline
	\cellcolor{egelightblue}{\textit{\textbf{Mathe}}} & \thead{\footnotesize sehr gut\\erfüllt} & \thead{\footnotesize gut\\erfüllt} & \thead{\footnotesize teilweise\\erfüllt} & \thead{\footnotesize nicht\\erfüllt} \\
	\hline
\end{xltabular}


%\begin{xltabular}{\textwidth}{|
%		>{\raggedright\arraybackslash\cellcolor{teal!70}\color{white}\bfseries}X
%		|Y|Y|Y|Y|}
%	\hline
%	Fach &
%	\theadcell{sehr gut\\erfüllt} &
%	\theadcell{gut\\erfüllt} &
%	\theadcell{teilweise\\erfüllt} &
%	\theadcell{nicht\\erfüllt} \\\hline%
%	%  …body of the table…
%	Deutsch      &   &   &   &   \\\hline%
%	Mathematik   &   &   &   &   \\\hline%
%\end{xltabular}



\newpagedefsLastPage
\thispagestyle{plain}

\fcolorbox{egeblue}{white}{%
	\begin{minipage}[t][.99\textheight][t]{.97\textwidth}
		%		% ---------- remarks ---------------------------------------
		\textbf{Bemerkungen:}\par\vspace{.5em}%
		\directlua{tex.print(tex_escape(student.comment))}\par
		
		\vfill
		
		%		% ---------- absence block ---------------------------------
		\begin{tabularx}{\linewidth}{lX}
			Versäumnisse: 
			& \directlua{tex.sprint(student.absenceDaysTotal)}~Tage (davon \directlua{tex.sprint(student.absenceDaysUnauthorized)}~Tage unentschuldigt)\\[.3em]
			& \directlua{tex.sprint(student.absenceHoursTotal)}~Stunden (davon \directlua{tex.sprint(student.absenceHoursUnauthorized)}~Stunden unentschuldigt)\\%
		\end{tabularx}
		\vspace{3cm}\par
		Erfurt, den \schoolReportDay
		
		%		% ---------- signatures 1 ---------------------------------
		\vspace{2.5cm}\par
		\begin{tabularx}{\linewidth}{C{.4\linewidth}Y C{.4\linewidth}}
			\hrulefill & & \hrulefill \\[6pt]
			Klassenleiter/in & & Klassenleiter/in\\
		\end{tabularx}
		%		% ---------- signatures 2 ---------------------------------
		\vspace{2.5cm}\par
		\begin{tabularx}{\linewidth}{C{.4\linewidth}Y C{.4\linewidth}}
			\rule{\linewidth}{1pt} & & \\[-.7ex]
			{\centering SchulleiterIn} & {\tiny Siegel} & \\%
		\end{tabularx}		
		%		% ---------- signatures 3 ---------------------------------
		\vspace{2.5cm}\par
		\begin{tabularx}{\linewidth}{C{.4\linewidth}X C{.4\linewidth}}
			\rule{\linewidth}{1pt} & & \rule{\linewidth}{1pt}\\%
			{\centering Erziehungsberechtigte} & &	{\centering SchülerIn}\\%
		\end{tabularx}		
		\vfill		
		%		% ---------- legal notice ---------------------------------
		{\footnotesize
			\textbf{Rechtsbehelf:}\\
			Gegen die Versetzungsentscheidung kann innerhalb eines Monats nach
			Bekanntgabe des Zeugnisses Widerspruch erhoben werden. Der Widerspruch
			ist bei der Evangelischen Gemeinschaftsschule Erfurt,
			Eugen-Richter-Straße 22, 99085 Erfurt schriftlich oder zur Niederschrift
			zu erheben.\\[3em]
			\begin{tabularx}{\linewidth}{C{.6\linewidth}}
				\rule{\linewidth}{1pt} \\[-.3ex]
				{\centering Kenntnisnahme der Sorgeberechtigten}\\%
			\end{tabularx}
		}
	\end{minipage}%
}%

	
\end{document}